\chapter{Literature Review}

The following chapter provides an in-depth review of foundational concepts and recent technological advances underpinning the research. The review aims to establish the theoretical base, examine pertinent frameworks and tools, and identify critical gaps justifying the present work.

\section{Theoretical Background}

\subsection{Agentic AI and Multi-Agent Systems (MAS)}

Artificial Intelligence agents are autonomous entities possessing environment perception, reasoning capabilities, and decision-making skills to pursue goals \cite{sycara1998multiagent}. Multi-Agent Systems (MAS) embody collections of such agents collaborating to solve complex, distributed tasks beyond single-agent capacity.

The Contract Net Protocol \cite{smith1980contract} remains a seminal contribution by formalizing decentralized task negotiation, allocation, and coordination among agents via bidding and task contracting, thereby laying the groundwork for distributed problem-solving architectures. The described negotiation model informs the collaborative agent interactions in the thesis, particularly in architect supervisor and domain-specific generator/validator agents.

\subsection{Retrieval-Augmented Generation (RAG)}

Large Language Models (LLMs) excel in understanding and generating human-like language but suffer from static training knowledge, causing hallucinations and inaccuracies when encountering novel or domain-specific queries.

Retrieval-Augmented Generation (RAG) \cite{lewis2020rag} mitigates such limitations by combining LLM capabilities with external knowledge sources. Upon receiving a query, the RAG framework retrieves relevant documents from curated knowledge bases, then conditions the language model's generation on the fresh, authoritative information to enhance factual correctness.

Recent enhancements like Self-RAG \cite{asai2023selfrag} introduce self-reflective mechanisms, allowing models to critique and iteratively improve their own outputs. Self-correction capabilities strongly support the iterative Evaluator-Optimizer strategy employed in the project.

\subsection{Reasoning and Acting (ReAct)}

Complex cognition requires coupling reasoning with action. The ReAct framework \cite{yao2023react} formalizes the cognition-action duality by enabling LLMs to generate explicit reasoning steps ('thoughts') interleaved with concrete actions such as tool calls or queries. The resulting cyclical reasoning-acting loop allows agents to refine results through observations and additional reasoning, forming the basis for generator and validator agents in the system.

\section{Implementation Frameworks}

\subsection{LangChain and LCEL}

LangChain \cite{langchain2024docs} is an extensible architecture providing building blocks for LLM applications, including prompt templates, model integration, and tool connectors. Its modern expression language, LCEL \cite{langchain2024lcel}, allows declarative specification of complex, pipelined workflows by composing modular components, enabling scalable parallel and sequential processes with ease.

\subsection{LangGraph: Enabling Agentic Loops}

While LCEL excels in linear workflows, many AI systems require iterative, stateful cycles. LangGraph \cite{langgraph2024docs} extends LangChain to support cyclic `StateGraph`s, where nodes represent agents modifying shared state, and edges encode communication and control flow—including conditional looping edges to implement iterative processes.

The core multi-agent feedback workflow (Evaluator-Optimizer loop) leverages LangGraph to allow repeated cycles of generation, validation, and refinement until convergence.

\section{Knowledge Base Construction}

\subsection{Web Scraping}

The accuracy and reliability of RAG-driven LLM workflows depend directly on the quality of the supporting knowledge base. The current work employs both automated and manual web scraping methods to continuously aggregate vendor-maintained documentation, whitepapers, and best practice guides from AWS, Azure, and potentially other providers, ensuring an up to date corpus.

\subsection{Docling: AI Document Parsing}

Complex vendor documents often appear in multi-format layouts (PDF, DOCX, PPTX), posing challenges for conventional parsers. Docling \cite{docling2025paper} is an AI-powered, open source document conversion tool that incorporates state-of-the-art models to parse such complex documents, extracting structured text, tables, and images accurately.

In the proposed framework, Docling integrates as a critical preprocessing step to generate clean, semantically meaningful document fragments suitable for downstream embedding and indexing \cite{docling2024web}.

\section{Local Model Development Environment}

\subsection{Ollama}

Ollama \cite{ollama2024web} streamlines the deployment of open-source LLMs such as Llama 3 and embedding-focused models like Nomic-Embed. The platform provides accessible APIs and abstracts hardware complexity, enabling efficient local embedding generation fundamental to RAG vector indices.

\subsection{ChromaDB: Vector Embedding Storage and Retrieval}

The framework utilizes the \textbf{ChromaDB} vector database \cite{chromadb2024} to store vector embeddings efficiently and facilitate fast similarity searches essential for RAG.

ChromaDB is a purpose-built, open-source vector database optimized for machine learning workloads, allowing scalable and low-latency retrieval of relevant documents based on approximate nearest neighbor searches. By indexing the Ollama-generated embeddings (derived from the \texttt{nomic-embed-text} model), ChromaDB enables Validator agents to retrieve contextually pertinent knowledge fragments rapidly, grounding assessments with high precision.

The integration of ChromaDB in the knowledge retrieval pipeline is critical for system performance and accuracy. The database supports complex queries where exact keyword matching fails, leveraging semantic similarity to connect diverse phrasings and domain-specific terminology prevalent in cloud service descriptions.

Thus, ChromaDB serves as the backbone vector store enabling reliable, validated factual grounding and iterative self-correction within the Evaluator-Optimizer architecture deployed.

\subsection{Apple Metal Performance Shaders (MPS)}

To achieve high throughput on commodity hardware without relying on costly GPUs, the system utilizes Apple’s Metal Performance Shaders (MPS) framework \cite{apple2024mps}. MPS allows ML workloads, particularly PyTorch-based embedding model computations, to run natively and efficiently on Apple Silicon GPUs, optimizing local development cycles.

\section{Review of Existing Work}

\subsection{Advanced RAG and Self-Correction}

Building on the original RAG framework \cite{lewis2020rag}, recent literature presents improvements emphasizing self-correction via reflection \cite{asai2023selfrag}. Such methods position self-evolving LLMs that critique and refine outputs, directly motivating the Evaluator-Optimizer approach to architecture recommendation in the thesis \cite{anthropic2024agents}.

\subsection{Multi-Agent Frameworks for LLMs}

Frameworks such as AutoGen \cite{wu2023autogen} enable multi-agent conversational workflows using LLMs, providing flexible agent orchestration but typically without explicitly defined cyclic state management required for iterative validation loops. The proposed architecture advances such designs by employing LangGraph \cite{langgraph2024docs} for granular and stateful multi-agent control.

\subsection{AI in Cloud Environment Management}

Recent proprietary solutions including Amazon Q \cite{aws2023q} and Microsoft Copilot for Azure \cite{microsoft2024copilot} exemplify AI integration in cloud ecosystems. Despite the inherent capabilities, proprietary tools operate within single cloud ecosystems and lack comprehensive multi-cloud solutions or architectural debate mechanisms—a gap the present research aims to fill.

\section{Identifying the Research Gap}

The conducted literature synthesis uncovers critical unmet needs:

\begin{enumerate}
    \item \textbf{Lack of Comprehensive Architectural Tools:} Current AI assistants provide vendor-specific, fragmented guidance lacking end-to-end validated cloud solution synthesis \cite{aws2023q, microsoft2024copilot}.
    \item \textbf{Absence of Dynamic, Argumentative Multi-Cloud Evaluation:} Earlier cloud brokering systems \cite{chauhan2014brokering} do not exhibit intelligent AI powered, iterative argumentative comparison capabilities crucial to unbiased multi-cloud decision-making.
    \item \textbf{Potential for Reliability via Advanced RAG and Self-Reflection:} Emerging methods \cite{lewis2020rag, asai2023selfrag} empower factual grounding and iterative self-correction, validating the Evaluator-Optimizer cyclical feedback methodology \cite{anthropic2024agents}.
\end{enumerate}

The present research thus targets advancing state-of-the-art by architecting an agentic multi-agent AI framework integrating cyclical validation, cross-provider architectural competition, and grounded knowledge retrieval to support superior decision-making in cloud architecture.